% GENERATED FILE. Edit canonical lessons, then run npm run generate:books.
% canonical-source-hash: fnv1a64:4dbdffa6914f74bf
% canonical-lessons: ES-C747-cunado, ES-C747-suegro, ES-C747-suegra, ES-C747-yerno, ES-C747-nuera

\chapter{La familia política}
\label{ch:es-la-familia-politica}
\hlchaptermodality{747}

\begin{chapteropening}
\textbf{By the end of this chapter:} \emph{I can name my in-laws and talk about them.}

Complete the last lesson of chapter 747: I can name my in-laws and talk about them.
\end{chapteropening}

\section[el cuñado]{el cuñado — the brother-in-law}

\label{lesson:ES-C747-cunado}



\begin{quote}
Before the new one: say the Spanish for to flirt, then the Spanish for the lover.
\end{quote}

\subsection*{You'll want to know: el cuñado}
\textbf{el cuñado} — \textquotedblleft{}the brother-in-law\textquotedblright{}.

\begin{quote}
Mi cuñado vive en Valencia.
\end{quote}

\begin{quote}
\emph{My brother-in-law lives in Valencia.}
\end{quote}

\emph{La cuñada} — sister-in-law.

\begin{cousinweb}
\emph{cuñado} is usually traced to Latin \emph{cognātus} \textquotedblleft{}related by blood\textquotedblright{} (Lewis \& Short); Spanish narrowed it to the in-law.
\end{cousinweb}

\begin{grammarlens}[title={What you've built}]
One more word for telling, explaining and describing at B1.
\end{grammarlens}

\subsection*{Your turn}
\begin{itemize}
\raggedright
  \item \emph{Say it:} \emph{el cuñado}
  \item \emph{Say it:} \emph{el cuñado}, once more
  \item \emph{Say it:} a sentence of your own with \emph{el cuñado}
  \item \emph{Recall:} say the Spanish for to flirt, then the Spanish for the lover, then say \emph{el cuñado} again
\end{itemize}

\begin{tcolorbox}[breakable,colback=teal!4,colframe=teal!35!black,title={Before you move on}]
What does el cuñado mean? (\textquotedblleft{}The brother-in-law\textquotedblright{}.) Say it once more.
\end{tcolorbox}
\section[el suegro]{el suegro — the father-in-law}

\label{lesson:ES-C747-suegro}



\begin{quote}
Before the new one: say the Spanish for the lover, then the Spanish for the brother-in-law.
\end{quote}

\subsection*{You'll want to know: el suegro}
\textbf{el suegro} — \textquotedblleft{}the father-in-law\textquotedblright{}.

\begin{quote}
Mi suegro cocina muy bien.
\end{quote}

\begin{quote}
\emph{My father-in-law cooks very well.}
\end{quote}

\emph{Los suegros} — the parents-in-law.

\begin{cousinweb}
\emph{suegro} is usually traced to Latin \emph{socer}, \emph{socerum} \textquotedblleft{}father-in-law\textquotedblright{} (Lewis \& Short); the stressed \emph{o} became \emph{ue}.
\end{cousinweb}

\begin{grammarlens}[title={What you've built}]
One more word for telling, explaining and describing at B1.
\end{grammarlens}

\subsection*{Your turn}
\begin{itemize}
\raggedright
  \item \emph{Say it:} \emph{el suegro}
  \item \emph{Say it:} \emph{el suegro}, once more
  \item \emph{Say it:} a sentence of your own with \emph{el suegro}
  \item \emph{Recall:} say the Spanish for the lover, then the Spanish for the brother-in-law, then say \emph{el suegro} again
\end{itemize}

\begin{tcolorbox}[breakable,colback=teal!4,colframe=teal!35!black,title={Before you move on}]
What does el suegro mean? (\textquotedblleft{}The father-in-law\textquotedblright{}.) Say it once more.
\end{tcolorbox}
\section[la suegra]{la suegra — the mother-in-law}

\label{lesson:ES-C747-suegra}



\begin{quote}
Before the new one: say the Spanish for the brother-in-law, then the Spanish for the father-in-law.
\end{quote}

\subsection*{You'll want to know: la suegra}
\textbf{la suegra} — \textquotedblleft{}the mother-in-law\textquotedblright{}.

\begin{quote}
Mi suegra viene a comer el domingo.
\end{quote}

\begin{quote}
\emph{My mother-in-law is coming for lunch on Sunday.}
\end{quote}

\begin{cousinweb}
\emph{suegra} is inherited from spoken Latin \emph{socra}, Latin \emph{socrus} \textquotedblleft{}mother-in-law\textquotedblright{} (EtymDB).
\end{cousinweb}

\begin{grammarlens}[title={What you've built}]
One more word for telling, explaining and describing at B1.
\end{grammarlens}

\subsection*{Your turn}
\begin{itemize}
\raggedright
  \item \emph{Say it:} \emph{la suegra}
  \item \emph{Say it:} \emph{la suegra}, once more
  \item \emph{Say it:} a sentence of your own with \emph{la suegra}
  \item \emph{Recall:} say the Spanish for the brother-in-law, then the Spanish for the father-in-law, then say \emph{la suegra} again
\end{itemize}

\begin{tcolorbox}[breakable,colback=teal!4,colframe=teal!35!black,title={Before you move on}]
What does la suegra mean? (\textquotedblleft{}The mother-in-law\textquotedblright{}.) Say it once more.
\end{tcolorbox}
\section[el yerno]{el yerno — the son-in-law}

\label{lesson:ES-C747-yerno}



\begin{quote}
Before the new one: say the Spanish for the father-in-law, then the Spanish for the mother-in-law.
\end{quote}

\subsection*{You'll want to know: el yerno}
\textbf{el yerno} — \textquotedblleft{}the son-in-law\textquotedblright{}.

\begin{quote}
Su yerno es médico.
\end{quote}

\begin{quote}
\emph{Her son-in-law is a doctor.}
\end{quote}

\begin{cousinweb}
\emph{yerno} is inherited from Latin \emph{gener}, \emph{generum} \textquotedblleft{}son-in-law\textquotedblright{} (EtymDB); Latin \emph{ge-} became \emph{ye-}.
\end{cousinweb}

\begin{grammarlens}[title={What you've built}]
One more word for telling, explaining and describing at B1.
\end{grammarlens}

\subsection*{Your turn}
\begin{itemize}
\raggedright
  \item \emph{Say it:} \emph{el yerno}
  \item \emph{Say it:} \emph{el yerno}, once more
  \item \emph{Say it:} a sentence of your own with \emph{el yerno}
  \item \emph{Recall:} say the Spanish for the father-in-law, then the Spanish for the mother-in-law, then say \emph{el yerno} again
\end{itemize}

\begin{tcolorbox}[breakable,colback=teal!4,colframe=teal!35!black,title={Before you move on}]
What does el yerno mean? (\textquotedblleft{}The son-in-law\textquotedblright{}.) Say it once more.
\end{tcolorbox}
\section[la nuera]{la nuera — the daughter-in-law}

\label{lesson:ES-C747-nuera}



\begin{quote}
Before the new one: say the Spanish for the mother-in-law, then the Spanish for the son-in-law.
\end{quote}

\subsection*{You'll want to know: la nuera}
\textbf{la nuera} — \textquotedblleft{}the daughter-in-law\textquotedblright{}.

\begin{quote}
Se lleva muy bien con su nuera.
\end{quote}

\begin{quote}
\emph{She gets on very well with her daughter-in-law.}
\end{quote}

\begin{cousinweb}
\emph{nuera} is inherited from Latin \emph{nurus} \textquotedblleft{}daughter-in-law\textquotedblright{}, through spoken \emph{nora} and Old Spanish \emph{nuera} (EtymDB).
\end{cousinweb}

\begin{grammarlens}[title={What you've built}]
One more word for telling, explaining and describing at B1.
\end{grammarlens}

\subsection*{Your turn}
\begin{itemize}
\raggedright
  \item \emph{Say it:} \emph{la nuera}
  \item \emph{Say it:} \emph{la nuera}, once more
  \item \emph{Say it:} a sentence of your own with \emph{la nuera}
  \item \emph{Recall:} say the Spanish for the mother-in-law, then the Spanish for the son-in-law, then say \emph{la nuera} again
\end{itemize}

\begin{tcolorbox}[breakable,colback=teal!4,colframe=teal!35!black,title={Before you move on}]
What does la nuera mean? (\textquotedblleft{}The daughter-in-law\textquotedblright{}.) Say it once more.
\end{tcolorbox}
